\documentclass[11pt]{article}

\usepackage[margin=1in]{geometry}
\usepackage[T1]{fontenc}
\usepackage{booktabs}
\usepackage{graphicx}
\usepackage{xcolor}
\usepackage{amsmath}
\usepackage{amssymb}
\usepackage{caption}
\usepackage{microtype}
\usepackage{url}
\usepackage{pgfplots}
\pgfplotsset{compat=1.18}
\usepackage[hidelinks]{hyperref}

\newcommand{\sys}{CS~Picks}
\newcommand{\csr}{CSRankings}

% Shared look for every plot, so the figures read as one family.
\pgfplotsset{
  cspicks/.style={
    width=0.47\textwidth, height=5.2cm,
    tick align=outside, tick pos=left,
    grid=major, grid style={gray!20},
    label style={font=\small}, tick label style={font=\footnotesize},
    title style={font=\small\bfseries},
    legend style={font=\footnotesize, draw=none, fill=none}
  }
}

\title{\sys: An Uncertainty-Aware Interface for Exploring\\
Computer Science Departments, Faculty, and Venues}

\author{ThanhVu Nguyen\\
George Mason University\\
\texttt{tvn@gmu.edu}}

\date{\today}

\begin{document}
\maketitle

\begin{abstract}
Metrics-based rankings of computer science departments have become a routine
input to consequential decisions, most visibly for prospective doctoral
students choosing where to apply. \csr{} is the field's de facto standard: it
is transparent, reproducible, and deliberately narrow, computing a single
ordered list from fractional publication credit at a curated set of venues.
Its very legibility, however, encourages readers to treat one integer as a
measurement of departmental quality.

We present \sys{}, a static, entirely client-side web application that reuses
\csr{}'s inputs and scoring rules but reorganizes them around the questions an
applicant actually asks. \sys{} makes the discretionary parameters of a ranking
manipulable and visible, and adds views that a single ordered list cannot
express: a \emph{rank stability} sweep that recomputes a department's position
under every combination of look-back window and venue set; per-subfield
attribution of both departmental score and individual faculty output; a
\emph{per-faculty} ordering that measures output per head rather than in
aggregate; a \emph{fragility} analysis reporting how many departures would move
a department out of its rank band; head-to-head comparison of two departments,
two researchers, or two subfields; and a funding view over NSF awards
attributed to the same roster.

Applying the stability sweep to all 190 ranked U.S.\ departments produces our
central empirical finding: rank stability is highly unevenly distributed. The
median department moves 17 positions across 16 defensible settings, and 77.4\%
move more than 10, yet the top-10 departments move a median of only 5. The
ordering is therefore trustworthy exactly where readers least need guidance and
unreliable across the broad middle where most applicants are choosing. We
report companion analyses at three levels of granularity --- department,
researcher, and subfield --- including a per-faculty ordering that disagrees
with the \csr{} ordering by a median of 13 positions, a subfield distribution
in which machine learning alone accounts for 24.0\% of all U.S.\ output, an
individual distribution in which the top 10\% of researchers produce 40.0\% of
it, and an NSF funding ordering that agrees with the publication ordering only
loosely ($\tau = 0.778$, median displacement 12 positions).
\end{abstract}

\section{Introduction}
\label{sec:intro}

A prospective doctoral student in computer science faces a decision with a long
horizon, poor feedback, and very little trustworthy information. Departmental
rankings fill that vacuum. They are compact, ordinal, and easy to act on, and
so they are acted on --- by applicants, by administrators, and by the
departments that manage themselves against
them~\cite{usnews,espeland2016engines}.

Within computer science, \csr{}~\cite{csrankings} occupies a distinctive
position. Unlike reputation surveys, it is computed from a public, auditable
procedure over public data: publications recorded in DBLP~\cite{dblp} at a
curated list of selective conferences, credited fractionally among coauthors,
aggregated per area, and combined across areas by a geometric mean. Anyone can
reproduce the number. This transparency is a genuine advance, and we build
directly on it.

Transparency of \emph{procedure}, however, is not the same as transparency of
\emph{contingency}. A reproducible pipeline still rests on discretionary
choices --- which venues count, how many years to look back, whether credit
follows an author's current or historical affiliation --- and the output
format, a single ordered list, cannot express that any of those choices were
made. A reader who sees a department at position 32 has no way to learn from
the artifact itself that the same department sits at 23 under one defensible
setting and 50 under another. The number's precision is real; its stability is
not, and nothing on the page distinguishes the two.

A single departmental ordering also answers only one of the three questions an
applicant is actually asking. Applicants join research groups, work with
individual advisors, and commit to a subfield. \sys{} therefore reports at all
three levels --- department, subfield, and individual researcher --- and lets
any two entities of the same kind be placed side by side.

This paper describes \sys{}, a companion tool to a guide for prospective
computer science doctoral students~\cite{phddemystify}. \sys{} is not an
alternative ranking. It consumes \csr{}'s inputs and reimplements its scoring
faithfully, then asks what else can be built once that computation is fast
enough to run many times in a browser. Our contributions are:

\begin{itemize}
  \item \textbf{A rank stability sweep} (\S\ref{sec:stability}) that recomputes
    a department's rank under all 16 combinations of four look-back windows and
    four venue sets, holding region fixed, and reports the resulting envelope
    rather than a point estimate.
  \item \textbf{Per-subfield attribution} (\S\ref{sec:attribution}) of both
    departmental score composition and individual faculty output, so that a
    name listed under a subfield reports what that person published
    \emph{in that subfield} rather than their department-wide total.
  \item \textbf{A per-faculty ordering} (\S\ref{sec:percapita}) that divides
    adjusted output by publishing headcount, separating how much a department
    produces from how many people it employs.
  \item \textbf{A fragility analysis} (\S\ref{sec:fragility}) reporting the
    number of departures that would move a department out of the top 10, 25, or
    50, computed without naming the individuals concerned.
  \item \textbf{Head-to-head comparison} (\S\ref{sec:comparemethod}) of two
    departments, two researchers, or two subfields, reporting per-area margins
    and, for departments, a decomposition of the rank gap into the areas that
    produced it.
  \item \textbf{A funding view} (\S\ref{sec:nsfmethod}) over NSF awards
    attributed to the same roster, which lets a publication-based ordering be
    checked against an independently produced one.
  \item \textbf{An empirical characterization} (\S\ref{sec:results}) of rank
    stability, venue sensitivity, output concentration, per-faculty ordering,
    fragility, the individual-researcher distribution, the subfield
    distribution, and funding--publication divergence across the U.S.\
    departments in the current data.
  \item \textbf{A fully client-side implementation} (\S\ref{sec:implementation})
    with no backend, in which a complete re-ranking of 20{,}822 faculty takes
    roughly 100\,ms, making interactive sensitivity analysis practical.
\end{itemize}

\section{Background}
\label{sec:background}

\subsection{\csr{} and its scoring rule}

\csr{} ranks departments by publication output at a curated set of highly
selective conferences, deliberately excluding journals, workshops, and citation
counts. Three design decisions matter for everything that follows.

\paragraph{Fractional credit.} A paper with $k$ authors contributes $1/k$ to
each author, and thence to their department. This is a deliberate
counterweight to the very different coauthorship norms across subfields, where
raw paper counts would reward large-team areas. We refer to the resulting
quantity as \emph{adjusted count} throughout.

\paragraph{Geometric mean across areas.} A department's score is
\begin{equation}
  S(d) = \left( \prod_{a \in A} \left( c_{d,a} + 1 \right) \right)^{1/|A|},
  \label{eq:score}
\end{equation}
where $A$ is the set of top-level research areas and $c_{d,a}$ is department
$d$'s adjusted count in area $a$. The product runs over \emph{all} areas, not
only those where the department is active, so a zero contributes a factor of
one. This strongly rewards breadth: a department strong in every area beats one
that is dominant in a few. It is the single most consequential choice in the
formula and is often misread as an average of strengths.
\S\ref{sec:comparison} shows it operating on a concrete pair of departments
whose total output differs by 0.3\%.

\paragraph{Venue curation.} Which conferences count is a maintained,
contested list. \csr{} distinguishes a default set from a larger set including
additional venues; independent efforts such as the CORE
rankings~\cite{coreportal} draw the line elsewhere. \sys{} supports four sets
(\S\ref{sec:venuesets}).

An important structural consequence of Eq.~\ref{eq:score} is that scores are
rounded to one decimal place before ranking, and ties share a rank under
standard competition ranking. In the tail of the distribution, where
departments differ by fractions of a publication, this rounding produces
substantial tied blocks --- a fact that becomes visible in our stability
analysis.

\subsection{Why a single ordered list under-serves the applicant}

The critiques of academic ranking are well rehearsed: rankings reshape the
behavior they claim to measure~\cite{goodhart}, and institutions reorganize
around them~\cite{usnews}. Our concern is narrower and more practical. Even
granting that adjusted publication count at selective venues is a reasonable
proxy for research activity, the \emph{presentation} of that proxy as a single
integer discards three things an applicant needs:

\begin{enumerate}
  \item \textbf{Sensitivity.} How much of position 32 is the department, and
    how much is the window and venue set?
  \item \textbf{Locality.} An applicant does not join a department; they join a
    research group. Departmental aggregates say little about whether one
    subfield is strong or whether one professor is taking students.
  \item \textbf{Scale.} A sum over faculty conflates a productive department
    with a large one, and says nothing about how much of the total rests on how
    few people.
\end{enumerate}

\sys{} addresses each in turn.

\section{Methods}
\label{sec:methods}

\subsection{Data sources and joins}

\sys{} fetches three \csr{} CSV files at page load --- the faculty roster
(name, affiliation, homepage, Scholar and ORCID identifiers), the per-author
publication table (name, area, year, raw count, adjusted count), and the
institution table (region, country, homepage) --- and joins them on author
name. Table~\ref{tab:dataset} summarizes the resulting corpus.

\begin{table}[t]
\centering
\caption{The corpus as fetched at the time of writing. Publication records are
per author, area, and year; they are aggregates, not individual papers. The
final block covers the separately synchronized NSF award snapshot
(\S\ref{sec:nsfmethod}).}
\label{tab:dataset}
\begin{tabular}{lr}
\toprule
Quantity & Value \\
\midrule
Faculty in roster            & 20{,}822 \\
Institutions                 & 702 \\
Publication records          & 296{,}267 \\
Year range                   & 1970--2026 \\
Faculty with a homepage      & 20{,}822 (100.0\%) \\
Faculty with a Scholar ID    & 18{,}256 (87.7\%) \\
\midrule
NSF awards in snapshot       & 22{,}696 \\
\quad within the 2016--2026 window   & 10{,}226 \\
Faculty matched to an award  & 4{,}101 \\
Institutions matched         & 182 \\
Attributed award value       & \$4.12B \\
\bottomrule
\end{tabular}
\end{table}

A consequence of the third row deserves emphasis, because it constrains every
method below: the publication table is \emph{aggregated per author, area, and
year}. It contains no paper titles, no paper identifiers, and no author order.
Any analysis requiring per-paper structure --- coauthorship, first-versus-last
authorship, cross-institutional collaboration --- cannot be answered from
\csr{} data alone and must go to DBLP directly, one author at a time.

Two optional sources extend the corpus. Affiliation histories derived from
OpenAlex~\cite{openalex} allow publications to be credited to the institution
an author held \emph{at publication time} rather than their current one; these
are generated offline and loaded on demand. NSF award
records~\cite{nsfawards}, synchronized offline and attributed fractionally
among principal investigators, support a separate funding view
(\S\ref{sec:nsfmethod}).

\subsection{The filtering and ranking pipeline}

The core computation, invoked by nearly every view, is a three-stage pipeline
parameterized by a year range $[y_0, y_1]$, a region $r$, a venue set $V$, and
optionally an affiliation history map $H$:

\begin{enumerate}
  \item \textbf{Collect.} Retain publications with $y_0 \le \text{year} \le
    y_1$ whose venue is in $V$; credit each to the institution(s) implied by
    $H$ (or the current affiliation when $H$ is absent), keeping only those in
    $r$. Accumulate per-department, per-area, and per-faculty totals in the
    same pass.
  \item \textbf{Score.} Apply Eq.~\ref{eq:score} to each department.
  \item \textbf{Rank.} Sort by score and assign standard competition ranks,
    overall and within each area, so tied scores share a rank and the next
    distinct score skips ahead.
\end{enumerate}

Crucially, per-faculty and per-area-per-faculty totals are accumulated
\emph{during} stage 1 rather than derived afterwards. Under historical
attribution a single author's output may be split across several institutions,
and that split cannot be reconstructed from their aggregate area totals once
the pass has finished. The same accumulation is what makes the fragility
analysis (\S\ref{sec:fragility}) and the per-subfield researcher tables
(\S\ref{sec:researchers}) computable at all.

\subsection{Venue sets}
\label{sec:venuesets}

\sys{} supports four venue sets: \csr{}'s default list; \csr{}'s extended list
including additional venues; CORE~A$^*$ only; and CORE~A together with
A$^*$. A venue that upstream scrapes but assigns to no research area is
counted by none of them, matching \csr{}'s own behavior.

\subsection{Rank stability}
\label{sec:stability}

Let $W = \{5, 10, 20, 30\}$ be look-back windows in years and $V$ the four
venue sets. For a fixed region $r$ and end year $y_1$, the sweep evaluates the
pipeline at every $(w, v) \in W \times V$, yielding 16 complete rankings. For a
department $d$ we report the envelope
\begin{equation}
  \big[\min_{(w,v)} \rho_{w,v}(d),\ \max_{(w,v)} \rho_{w,v}(d)\big],
\end{equation}
its median, and the \emph{spread} $\max - \min$, where $\rho_{w,v}(d)$ is $d$'s
rank in the corresponding run.

Two design decisions are worth stating explicitly.

\paragraph{Region is held fixed.} A rank among U.S.\ departments and a rank
worldwide answer different questions over different populations; pooling them
into a single interval would mix a \emph{scope} choice into what is meant to be
a range over \emph{methodological} choices. The user's selected region is
therefore held constant across the sweep.

\paragraph{Departments that vanish are reported, not dropped.} A department
active only at venues outside a given set does not appear in that run's ranking
at all. Silently omitting such runs would understate volatility, so we report
the count of settings under which a department is unranked alongside its
envelope.

One sweep produces ranks for every department simultaneously, so its cost is
amortized: the result is cached per region and reused for any department the
user subsequently inspects.

\subsection{Per-subfield attribution}
\label{sec:attribution}

Departmental cards list each active subfield with the faculty contributing to
it. Two quantities are distinguished that a naive implementation conflates.
First, a faculty member listed under subfield $a$ reports their output
\emph{within $a$ at that institution}, not their department-wide total ---
which, for a broad researcher, may differ by an order of magnitude. Second, we
report each subfield's share of the department's own score,
\begin{equation}
  \sigma_{d,a} = \frac{\ln(c_{d,a} + 1)}{\sum_{a' \in A} \ln(c_{d,a'} + 1)},
\end{equation}
which follows from taking logarithms of Eq.~\ref{eq:score}. Because these
shares are internal to one department and sum to one, they describe
composition, not standing: a large $\sigma_{d,a}$ means area $a$ drives $d$'s
score, not that $d$ leads others in $a$. The interface states this explicitly,
because the distinction is easy to lose.

\subsection{Per-faculty ordering}
\label{sec:percapita}

Eq.~\ref{eq:score} accumulates adjusted counts and then takes a geometric mean
across areas. Because the per-area counts are sums over faculty, a department
can raise its score by employing more people at a constant rate of output. The
resulting order therefore answers ``how much does this department publish'',
which is not the same question as ``how much does a member of this department
publish'' --- and for an applicant, who will work with a handful of people
rather than the whole department, the second is at least as relevant.

We therefore also order departments by
\begin{equation}
  P(d) = \frac{\sum_{a \in A} c_{d,a}}{|F_d|},
\end{equation}
where $F_d$ is the set of faculty at $d$ with at least one credited publication
in the selected window. Headcount is measured as \emph{publishing} faculty
rather than roster size, so that a department is not penalized for listing
members whose work falls outside the counted venues.

Departments with fewer than five publishing faculty are excluded from this
ordering rather than ranked. With two or three people the ratio is
substantially one person's output, and the resulting list is dominated by small
groups rather than by productive ones. This threshold is a judgement, and we
report it explicitly because it materially determines the head of the ordering.

\subsection{Fragility}
\label{sec:fragility}

Rank stability (\S\ref{sec:stability}) asks how much of a position is an
artifact of methodology. Fragility asks a complementary question about
personnel: how many departures would move this department out of its rank band?

Because removing faculty from one department leaves every other department's
score untouched, a hypothetical rank is a lookup rather than a re-ranking. We
proceed greedily. At each step we evaluate removing each remaining faculty
member, recompute the department's score under Eq.~\ref{eq:score} from its
per-area counts less that person's per-area contribution, and remove whichever
costs the department most. The choice is not simply the most prolific member:
under a geometric mean, losing the only contributor in a thin area can cost more
than losing a larger producer in a crowded one. We report the number of
departures needed to leave the top 10, 25, and 50, and the rank trajectory of
the first several steps.

Computing this requires per-person, per-area credit, which the pipeline
accumulates during stage 1 (\S\ref{sec:methods}) rather than deriving
afterwards.

\paragraph{Names are deliberately withheld.} The procedure necessarily
identifies which individuals carry a department, and the interface does not
display them. Publishing a ranked list of the people whose departure would most
damage their employer invites precisely the personnel conclusions this work
should not support, and the aggregate quantity --- how concentrated a
department's output is --- is what carries the analytic content. We report
distributions across departments in \S\ref{sec:results} for the same reason,
and \S\ref{sec:ethics} explains why the researcher tables of
\S\ref{sec:researchers} are treated differently.

\subsection{Researcher and subfield metrics}
\label{sec:researchermethod}

The same pipeline produces per-researcher and per-subfield aggregates, since
stage 1 accumulates both. For a researcher we report adjusted count, paper
count, \emph{breadth} (the number of top-level areas with any credit),
\emph{primary-area share} (the fraction of adjusted output in their largest
area), and a normalized-entropy \emph{balance} over the area vector. Two
temporal quantities are computed over the selected window: \emph{consistency},
the fraction of years with at least one credited publication, and
\emph{momentum}, the change in adjusted output between the trailing three-year
window and the three years before it. We call a researcher a \emph{specialist}
when primary-area share is at least 60\%, and a \emph{generalist} when breadth
is at least 4; the thresholds are reporting conventions, not claims.

For a subfield we report region-wide adjusted total, the number of departments
and distinct researchers with any credit, growth against the equal-length
preceding period, and the share of the subfield's national output held by its
five largest departments.

\subsection{Head-to-head comparison}
\label{sec:comparemethod}

Two entities of the same kind are compared by taking the union of their active
areas, ordering it by combined credit, and reporting the per-area margin
$c_{A,a} - c_{B,a}$. Alongside it we report a small scoreboard (rank, papers,
adjusted count, publishing faculty, and areas led, for departments; affiliation,
papers, adjusted count, active areas, and areas led, for researchers).

For departments we additionally decompose the score gap. Taking logarithms of
Eq.~\ref{eq:score}, the difference in scores is a sum of per-area terms
\begin{equation}
  \ln S(A) - \ln S(B) = \frac{1}{|A|}\sum_{a} \big[\ln(c_{A,a}+1) - \ln(c_{B,a}+1)\big],
  \label{eq:gap}
\end{equation}
so ranking areas by $|\ln(c_{A,a}+1) - \ln(c_{B,a}+1)|$ identifies exactly the
areas that produced the rank gap. This is worth stating because the answer is
frequently counterintuitive: the area where the raw margin is largest is
usually \emph{not} the area that moved the rank, since the logarithm compresses
differences at high volume and magnifies them near zero.
\S\ref{sec:comparison} gives a case where the two answers point in opposite
directions.

Subfields are compared by a separate procedure, reporting each side's
region-wide total, growth, participating departments and researchers, and the
researchers and departments active in both.

\subsection{NSF award attribution}
\label{sec:nsfmethod}

The funding view is built from a snapshot of the NSF award
database~\cite{nsfawards}, synchronized offline by querying the award-search
API once per (faculty, institution) pair over the \csr{} roster. Each award's
value is split equally among its listed investigators,
\begin{equation}
  \text{share}(p, w) = \frac{\text{value}(w)}{|\,\text{investigators}(w)\,|},
\end{equation}
and a person's attributed total is the sum of their shares. Value is the
estimated total award amount where present and the obligated amount otherwise.
An institution's attributed total is the sum over its matched faculty.

Equal splitting is a deliberate simplification and the main source of error in
this view: NSF publishes no per-investigator allocation, so a weighted split
would require inventing weights. We therefore report the equal share as a
\emph{share}, and separately report the full project value, which is the more
appropriate number for judging the scale of the work a group participates in.
The two differ substantially for large collaborative awards, and the interface
shows both.

A name-reconciliation step maps award investigators onto the current \csr{}
roster. \csr{} spells some faculty differently in its roster file than in its
publication table, and only names verified against the publication table are
retained, so an award with no resolvable investigator is dropped rather than
attributed to an approximate match.

\section{Implementation}
\label{sec:implementation}

\subsection{Architecture}

\sys{} is a static, multi-page web application with no backend of any kind. All
data is fetched client-side from public sources and processed in the browser;
deployment is a directory of static files. This is a substantive design choice
rather than an economy: it means the computation a reader sees is the
computation they can inspect, there is no server-side state that could
personalize or silently alter a ranking, and the tool cannot accumulate a
record of who looked up whom.

The application comprises four pages --- search and analysis, a ranking-impact
simulator, a discoveries feed, and an NSF funding explorer --- each with its
own entry module, sharing a common filter bar, chart layer, and autocomplete
component. The canonical data load is memoized, so modules loaded together
share a single download and parse.

\subsection{Performance}

A complete pass of the pipeline over all 20{,}822 faculty and 296{,}267
publication records --- filter, score, and rank --- takes approximately
100\,ms. This number is what makes the stability sweep viable: 16 full
re-rankings complete in about 1.7\,s.

Because that much uninterrupted computation would freeze the interface, the
sweep yields to the event loop between runs and reports incremental progress,
and its result is cached per region. The cost is therefore paid once per
region per session, and every department examined afterwards resolves from
cache.

\subsection{Offline generators and correctness}

Two datasets are too large or too slow to build at page load and are generated
offline: OpenAlex-derived affiliation histories, and the NSF award snapshot.

Correctness is maintained by a unit suite covering the filtering, scoring, and
ranking logic, rendering safety, and the metric definitions introduced here,
together with an end-to-end browser suite exercising the principal user flows.
Every number reported in \S\ref{sec:results} is emitted by a script that
imports the shipping modules and runs them against the live \csr{} inputs, so
the paper cannot drift from the implementation without the discrepancy being
visible on the next regeneration.

\section{Web interface}
\label{sec:ui}

\subsection{Search and filtering}

A single search box resolves institutions, faculty, research areas, and
individual venues, disambiguating by type. A shared filter bar exposes the
discretionary parameters --- region, year range, venue set, whether to show
rank ordinals, and whether to use historical affiliation --- and persists them
across pages and into the URL, so any view is shareable as a link that encodes
the settings that produced it. Making these parameters ambient rather than
buried is itself an argument: the reader is continuously reminded that the
numbers are conditional.

\subsection{Departmental cards}

A department expands to show, in order: the complete faculty roster ordered by
adjusted count, with each member's rank within the department when rank display
is enabled; each active subfield with its contributing faculty and their
\emph{within-subfield} output; and the subfield share of the department's own
score (\S\ref{sec:attribution}) rendered as a ranked bar chart.

The bars in that chart are scaled against the largest share rather than against
100\%. Because the shares derive from $\ln(c+1)$, the largest is typically
under 8\%, and drawing each bar at its literal percentage of the available
width renders two dozen visually indistinguishable stubs. This is a small
detail with a real consequence: the corrected scaling is what makes the
composition legible at a glance.

\subsection{Analysis panel}

Selecting a department or researcher opens a tabbed panel: activity,
faculty diversity, publishing effort, conference trends, area growth,
collaboration proxies, and rank stability, with the fragility analysis of
\S\ref{sec:fragility} beneath it. Tabs that do not apply to the current target
type are hidden rather than shown empty --- a researcher has no faculty
diversity --- and the researcher variant of the activity tab reports the
pattern metrics of \S\ref{sec:researchermethod} instead of a departmental rank
trajectory.

The default view, shown before any search, ranks departments and researchers in
two columns. A control above it switches the departmental column between the
\csr{} score and the per-faculty ordering of \S\ref{sec:percapita}; the
control appears only in that view, since a filtered result list is not a
ranking.

\subsection{Comparison}

A query of the form \texttt{A vs B} switches into a head-to-head view. The
autocomplete completes only the trailing side, so a comparison can be typed
without retyping the first entity, and the same syntax works for two
departments, two researchers, or two subfields; mixed-type comparisons are
refused with an explanation rather than silently coerced. The view renders the
scoreboard, the per-area margin chart, and --- for departments --- the rank-gap
decomposition of Eq.~\ref{eq:gap}. Worked examples appear in
\S\ref{sec:comparison}.

\subsection{Simulator and discoveries}

A separate simulator estimates how adding, removing, or transferring faculty
would change a department's rank, resolving hypothetical candidates against
DBLP so that real researchers not yet in \csr{} can be evaluated. Because
removing faculty leaves every other department's score fixed, the recomputation
is a rescore of the affected departments plus a re-sort, and runs
interactively.

A discoveries page surfaces notable movements without requiring the reader to
know what to search for: rank climbers and decliners, output and breadth
changes, subfield growth and consolidation, changes of subfield leadership, and
--- for U.S.\ and worldwide scope --- the funding sections of
\S\ref{sec:nsf}.

\subsection{Funding search}
\label{sec:fundingui}

The funding page searches the NSF snapshot independently of the publication
data; the search page never loads it. Its autocomplete offers three groups ---
universities, principal investigators, and NSF program names --- so
\texttt{Secure \& Trustworthy Cyberspace} is as valid a query as
\texttt{Carnegie Mellon University}. Results render as two parallel lists, one
of institutions and one of faculty, each showing award count, attributed share,
and full project value. The \texttt{A vs B} syntax works here too, over either
institutions or investigators.

A permanently visible data-health panel reports snapshot freshness, roster
coverage, and the matching caveats of \S\ref{sec:nsfmethod}, because a funding
figure invites more precise interpretation than it can support.

\section{Features and results}
\label{sec:results}

All figures in this section were computed by a script that runs the shipping
implementation against the live \csr{} inputs; none are estimated or
illustrative. Unless noted, the region is the United States, the end year is
2026, and the baseline setting is a 10-year window with \csr{}'s default venue
set. Institution names are abbreviated as the interface abbreviates them.

\subsection{Rank stability is unevenly distributed}
\label{sec:stabilityresults}

Table~\ref{tab:stability} summarizes the sweep over all 190 ranked U.S.\
departments.

\begin{table}[t]
\centering
\caption{Rank spread (worst minus best position) across 16 settings --- four
look-back windows $\times$ four venue sets --- for 190 ranked U.S.\
departments. Strata are defined by baseline rank.}
\label{tab:stability}
\begin{tabular}{lrr}
\toprule
Stratum & Departments & Median spread \\
\midrule
Baseline rank $\le$ 10   & 10  & 5.0 \\
Baseline rank $\le$ 25   & 26  & 6.0 \\
Baseline rank $\le$ 50   & 51  & 10.0 \\
All ranked departments   & 190 & 17.0 \\
\midrule
\multicolumn{3}{l}{\emph{Distribution over all departments}} \\
First quartile of spread   & \multicolumn{2}{r}{11 positions} \\
Median spread              & \multicolumn{2}{r}{17 positions} \\
Third quartile of spread   & \multicolumn{2}{r}{25 positions} \\
Maximum spread             & \multicolumn{2}{r}{46 positions} \\
Departments moving $>$10   & \multicolumn{2}{r}{77.4\%} \\
Departments moving $>$25   & \multicolumn{2}{r}{23.2\%} \\
\bottomrule
\end{tabular}
\end{table}

The headline number is the median spread of 17 positions, and the fact that
over three quarters of departments move by more than 10. But the stratification is
the more interesting result: the top of the ranking is markedly more stable
than the middle. Departments with a baseline rank of 10 or better move a median
of 5 positions; the field as a whole moves a median of 17.
Figure~\ref{fig:stability} shows the full picture --- the left panel gives the
cumulative distribution of spread, the right panel plots each department's
spread against its baseline rank, where the widening cone is the stratification
result rendered directly.

\begin{figure}[t]
\centering
\begin{tikzpicture}
\begin{axis}[cspicks, xlabel={Rank spread (positions)},
  ylabel={Departments at most (\%)}, ymin=0, ymax=100, xmin=0,
  title={Distribution of spread}]
\addplot[thick, blue!70!black] table [x=spread, y=share] {data/spread-cdf.dat};
\end{axis}
\end{tikzpicture}
\hfill
\begin{tikzpicture}
\begin{axis}[cspicks, xlabel={Baseline rank}, ylabel={Rank spread},
  xmin=0, ymin=0, title={Spread against baseline rank}]
\addplot[only marks, mark=*, mark size=1.1pt, blue!70!black, opacity=0.55]
  table [x=baseline, y=spread] {data/spread-by-rank.dat};
\end{axis}
\end{tikzpicture}
\caption{Rank spread across the 16-setting sweep, over 190 ranked U.S.\
departments. Stability is not a property of the ranking; it is a property of
where in the ranking a department sits.}
\label{fig:stability}
\end{figure}

The ordering is thus most trustworthy precisely where it is least decision-
relevant. An applicant weighing the top handful of departments is looking at a
comparatively robust ordering --- and is also choosing among options that a
single ranking position barely distinguishes. An applicant weighing options
ranked between 30 and 80, which describes the large majority of applicants, is
reading an ordering in which a 20-position difference can be produced entirely
by moving the window from 10 years to 20.

Even at the top, stability is not absolute: of the 10 departments in the
baseline top 10, only 5 remain in the top 10 under all 16 settings.

Table~\ref{tab:examples} gives individual envelopes.

\begin{table}[t]
\centering
\caption{Rank envelopes for selected departments across the 16 settings.}
\label{tab:examples}
\begin{tabular}{lrrrrr}
\toprule
Department & Baseline & Best & Worst & Median & Spread \\
\midrule
CMU          & 1  & 1  & 1  & 1  & 0 \\
UIUC         & 2  & 2  & 2  & 2  & 0 \\
Northeastern & 11 & 8  & 18 & 11 & 10 \\
GMU          & 32 & 23 & 50 & 42 & 27 \\
\bottomrule
\end{tabular}
\end{table}

CMU and UIUC are genuinely invariant: they hold their positions
under every setting tested. GMU ranges from 23rd to 50th depending on
nothing but the window and venue set, and its \emph{median} position across the
sweep is 42 --- ten places below the baseline value a reader would see. So the
baseline is not merely imprecise here; it is optimistic relative to the
methodology's own central tendency. Reporting only 32 conveys a precision the
data does not support, and this is the ordinary case, not a pathological one.

\subsection{Venue sets correlate globally but move individual departments}
\label{sec:venue}

Table~\ref{tab:venue} compares venue sets pairwise at a fixed 10-year window.

\begin{table}[t]
\centering
\caption{Pairwise agreement between venue sets at a 10-year window. $\tau$ is
Kendall's rank correlation~\cite{kendall1938}; overlap is the number of
departments common to both top-10 lists.}
\label{tab:venue}
\small
\begin{tabular}{llrrrr}
\toprule
Set A & Set B & $\tau$ & Median $|\Delta\text{rank}|$ & Max $|\Delta\text{rank}|$ & Top-10 overlap \\
\midrule
\csr{} default & \csr{} all    & 0.982 & 2 & 18 & 9 \\
\csr{} default & CORE A$^*$    & 0.953 & 5 & 20 & 9 \\
\csr{} default & CORE A/A$^*$  & 0.973 & 2 & 20 & 9 \\
\csr{} all     & CORE A$^*$    & 0.952 & 5 & 19 & 8 \\
\csr{} all     & CORE A/A$^*$  & 0.973 & 3 & 15 & 10 \\
CORE A$^*$     & CORE A/A$^*$  & 0.961 & 4 & 15 & 8 \\
\bottomrule
\end{tabular}
\end{table}

Read as aggregate correlations, these numbers are reassuring: $\tau$ never
falls below 0.95, so the venue sets broadly agree about the shape of the field.
Read as individual outcomes, they are not. Switching from \csr{}'s default set
to CORE~A$^*$ moves the median department 5 positions and some department 20
positions, and changes the membership of the top 10.

Table~\ref{tab:toptens} makes the last point concrete by printing all four
top-tens side by side. These are the same 190 departments, the same 10-year
window, and the same scoring formula; only the venue list differs.

\begin{table}[t]
\centering
\small
\caption{The U.S.\ top ten under each of the four venue sets, 10-year window.
Ties share a rank, so a column may list eleven names. Departments appearing in
some columns but not others are set in italics.}
\label{tab:toptens}
\begin{tabular}{rllll}
\toprule
& \csr{} default & \csr{} all & CORE A$^*$ & CORE A/A$^*$ \\
\midrule
1  & CMU          & CMU          & CMU              & CMU \\
2  & UIUC         & UIUC         & UIUC             & UIUC \\
3  & UCSD         & UCSD         & UCSD             & UCSD \\
4  & Georgia Tech & Georgia Tech & MIT              & Georgia Tech \\
5  & MIT          & UW \& MIT    & Georgia Tech     & MIT \\
6  & UW \& Berkeley & ---        & UMich            & UMD \\
7  & ---          & UMD          & UC Berkeley      & UMich \\
8  & UMD          & UC Berkeley  & \emph{Cornell} \& UW & UW \\
9  & UMich        & Northeastern & ---              & UC Berkeley \\
10 & Stanford     & UMich        & Stanford \& \emph{Purdue} & \emph{Northeastern} \\
\bottomrule
\end{tabular}
\end{table}

Four departments --- CMU, UIUC, UCSD, and Georgia Tech (except under CORE~A$^*$,
where it falls to 5) --- are effectively fixed. Below them the list is
unsettled. Stanford is 10th under the default set and under CORE~A$^*$ but
appears in neither of the other two top tens. Northeastern is 9th under
\csr{}'s full venue list and 10th under CORE~A/A$^*$, but is outside the top ten
under the default set and under CORE~A$^*$. Cornell and Purdue appear only
under CORE~A$^*$. A reader who has internalized ``the top ten'' as a stable
object is reading an artifact of one venue list.

\subsubsection{What \csr{} and CORE~A$^*$ actually disagree about}
\label{sec:venuecomposition}

The rank movement above has a cause that is visible in the venue lists
themselves, before any ranking is computed. Table~\ref{tab:venuecomp} sets the
two lists side by side per area.

\begin{table}[t]
\centering
\small
\setlength{\tabcolsep}{4pt}
\caption{Venue counts per area under each set, with the venues on which \csr{}
(either list) and CORE~A$^*$ disagree. Areas are ordered by how much the two
disagree; the five areas where they agree exactly are omitted. Venues marked
$\dagger$ are in \csr{}'s extended list but not its default one.}
\label{tab:venuecomp}
\begin{tabular}{lrrrr ll}
\toprule
& \multicolumn{4}{c}{Venue count} & \multicolumn{2}{c}{Venues in one list only} \\
\cmidrule(r){2-5}\cmidrule(l){6-7}
Area & Def. & All & A$^*$ & A/A$^*$ & \csr{} only & CORE A$^*$ only \\
\midrule
AI                    & 2 & 2 & 5 & 10 & ---                                          & aamas, kr, icaps \\
HPC                   & 3 & 3 & 0 & 3  & sc, hpdc, ics                                & --- \\
Mobile Computing      & 3 & 3 & 4 & 4  & mobisys                                      & percom, ipsn \\
Operating Systems     & 2 & 5 & 2 & 9  & eurosys$\dagger$, fast$\dagger$, usenixatc$\dagger$ & --- \\
Graphics              & 2 & 3 & 5 & 7  & eurographics$\dagger$                        & acmmm, ismar \\
Robotics              & 3 & 3 & 0 & 1  & icra, iros, rss                              & --- \\
Machine Learning      & 3 & 4 & 6 & 8  & ---                                          & colt, icdm \\
Design Automation     & 2 & 2 & 0 & 4  & dac, iccad                                   & --- \\
Embedded \& Real-Time & 3 & 3 & 1 & 2  & emsoft, rtas                                 & --- \\
Programming Languages & 2 & 4 & 2 & 6  & oopsla$\dagger$, icfp$\dagger$               & --- \\
Networks              & 2 & 2 & 2 & 3  & nsdi                                         & infocom \\
Cryptography          & 2 & 2 & 0 & 3  & crypto, eurocrypt                            & --- \\
Comp.\ Biology        & 2 & 2 & 0 & 1  & ismb, recomb                                 & --- \\
NLP                   & 3 & 3 & 2 & 5  & naacl                                        & --- \\
Architecture          & 3 & 4 & 3 & 4  & micro                                        & --- \\
Measurement           & 2 & 2 & 1 & 2  & imc                                          & --- \\
Software Engineering  & 2 & 4 & 3 & 15 & issta$\dagger$                               & --- \\
Algorithms            & 3 & 3 & 4 & 11 & ---                                          & podc \\
Econ \& Computation   & 2 & 2 & 1 & 1  & wine                                         & --- \\
HCI                   & 3 & 3 & 2 & 6  & ubicomp                                      & --- \\
Visualization         & 2 & 2 & 0 & 1  & vis                                          & --- \\
CS Education          & 1 & 1 & 0 & 4  & sigcse                                       & --- \\
\midrule
\textbf{Total venues} & \textbf{64} & \textbf{77} & \textbf{58} & \textbf{148} & 30 & 11 \\
\bottomrule
\end{tabular}
\end{table}

The headline numbers are close --- \csr{}'s full list carries 77 venues against
CORE~A$^*$'s 58 --- but the overlap is only 47. Thirty \csr{} venues are absent
from CORE~A$^*$ and eleven CORE~A$^*$ venues are absent from \csr{}, so the two
lists differ on 41 venues, over a third of their union. The two agree exactly
on only five of 27 areas: computer vision, information retrieval, security,
databases, and logic and verification.

Two structural differences do most of the work. First, CORE~A$^*$ is broader in
AI and theory --- it adds AAMAS, KR, and ICAPS to AI, COLT and ICDM to machine
learning, and PODC to algorithms --- which is why the AI-heavy departments in
Table~\ref{tab:toptens} hold their positions under it. Second, and far more
consequentially, \emph{CORE~A$^*$ has no venue at all in seven of the 27 areas}:
HPC, robotics, design automation, cryptography, computational biology,
visualization, and CS education. \csr{} counts 15 venues across those areas
(SC, ICRA, IROS, RSS, DAC, ICCAD, CRYPTO, EUROCRYPT, and others); CORE~A$^*$
counts none.

Under Eq.~\ref{eq:score} an area with no venues is not merely unmeasured, it is
inert: $c_{d,a} = 0$ for every department, so the factor $(c_{d,a}+1)$ equals 1
everywhere and the area stops discriminating between departments entirely.
Switching to CORE~A$^*$ therefore does not reweight seven areas, it deletes
them, and the geometric mean silently becomes an average over 20 effective
areas rather than 27.

This is measurable in exactly the way it should be. Across the 188 departments
ranked under both sets, the median department takes 16.8\% of its adjusted
output in those seven areas, and a department's exposure to them predicts how
far it falls: the correlation between dropped-area share and rank shift is
$r = 0.495$. The 72 departments with at least 20\% of their output in the
dropped areas move a median of 2 positions worse under CORE~A$^*$, while the 50
departments with under 10\% move a median of 7 positions better --- a
nine-position swing driven by nothing but which subfields a reader's chosen
venue list happens to cover.

For an applicant in robotics or cryptography the practical reading is blunt: a
CORE~A$^*$ ranking contains no information about their subfield whatsoever, and
a department's position under it reflects only the work it does elsewhere.
This is the strongest argument in the paper for exposing the venue set as a
visible control rather than a buried constant.

Table~\ref{tab:venuemovers} reports the departments that CORE~A$^*$ treats most
differently from the \csr{} default.

\begin{table}[t]
\centering
\small
\caption{Departments moved most by switching from \csr{}'s default venue set to
CORE~A$^*$, 10-year window. Negative shift means the department ranks better
under CORE~A$^*$.}
\label{tab:venuemovers}
\begin{tabular}{lrrrrr}
\toprule
Department & \csr{} default & \csr{} all & CORE A$^*$ & CORE A/A$^*$ & Shift \\
\midrule
\multicolumn{6}{l}{\emph{Gains under CORE~A$^*$}} \\
Auburn                & 116 & 113 & 101 & 105 & $-15$ \\
Brigham Young         & 116 & 113 & 101 & 115 & $-15$ \\
Cincinnati            & 116 & 120 & 101 & 115 & $-15$ \\
Drexel                & 89  & 89  & 75  & 84  & $-14$ \\
Emory                 & 83  & 82  & 71  & 75  & $-12$ \\
UC Santa Cruz         & 37  & 38  & 27  & 30  & $-10$ \\
\midrule
\multicolumn{6}{l}{\emph{Losses under CORE~A$^*$}} \\
Memphis               & 108 & 113 & 128 & 115 & $+20$ \\
Washington State      & 86  & 86  & 101 & 87  & $+15$ \\
Delaware              & 70  & 76  & 85  & 75  & $+15$ \\
Florida Atlantic      & 136 & 141 & 150 & 156 & $+14$ \\
Northern Arizona      & 136 & 141 & 150 & 138 & $+14$ \\
Alabama               & 136 & 141 & 150 & 138 & $+14$ \\
\bottomrule
\end{tabular}
\end{table}

Two things stand out. First, the largest movers cluster in the tail, where the
tied blocks noted in \S\ref{sec:background} mean a small score change crosses
many departments at once: the five departments tied at 136 under the default
set move together. Second, the movement is not confined to the tail. UC Santa
Cruz sits at 37 under the default set and 27 under CORE~A$^*$ --- a ten-place
move within the range applicants actually compare, produced entirely by
substituting one defensible venue list for another.

This is the familiar gap between a distribution and a case. A department
evaluating itself, or an applicant evaluating a department, experiences the
20-position move, not the 0.953 correlation. High rank correlation is
compatible with substantial individual displacement, and reporting only the
former is how sensitivity analyses come to reassure when they should caution.

\subsection{Output concentration rises sharply down the ranking}
\label{sec:concentration}

Table~\ref{tab:concentration} reports the share of a department's adjusted
output produced by its most prolific faculty, over 166 U.S.\ departments with
at least five active faculty.

\begin{table}[t]
\centering
\caption{Concentration of adjusted publication count, for 166 U.S.\
departments with at least 5 active faculty in a 10-year window.}
\label{tab:concentration}
\begin{tabular}{lr}
\toprule
Quantity & Value \\
\midrule
Median faculty per department          & 24 \\
Median top-1 share                     & 15.2\% \\
Median top-3 share                     & 36.7\% \\
Departments with top-3 share $>$ 50\%  & 31.9\% \\
\midrule
Median top-3 share, rank $\le$ 25      & 15.5\% \\
Median top-3 share, rank $>$ 100       & 60.3\% \\
\bottomrule
\end{tabular}
\end{table}

Nearly a third of departments derive more than half their adjusted output from
three people, and the contrast across the ranking is stark: a median top-3
share of 15.5\% among the top 25 versus 60.3\% below rank 100. Departments in
the tail are not scaled-down versions of departments at the top; they are
structurally different, with output concentrated in a few individuals.

For an applicant this is arguably more actionable than the rank itself. A
department ranked 120 whose standing rests on three faculty is a good choice if
one of those three takes the applicant as a student and a poor one otherwise,
and the rank alone cannot distinguish these outcomes. It also bears on the
stability result above: concentrated departments are precisely those whose
position is most sensitive to which venues are counted, since a few
individuals' venue preferences determine the department's profile.

\subsection{Measuring per faculty reorders the field}
\label{sec:percapitaresults}

Table~\ref{tab:percapita} places the two departmental orderings side by side
over the 166 U.S.\ departments with at least five publishing faculty.

\begin{table}[t]
\centering
\small
\caption{The two departmental orderings side by side, 10-year window, over the
166 U.S.\ departments with at least five publishing faculty. The left block
orders by the \csr{} score of Eq.~\ref{eq:score} and reports where each
department sits per faculty; the right block orders by adjusted output per
publishing faculty member and reports where each sits in the \csr{} order.}
\label{tab:percapita}
\begin{tabular}{r@{\hskip 1.5em}lrr@{\hskip 2.5em}lrrr}
\toprule
& \multicolumn{3}{c}{Ordered by \csr{} score} & \multicolumn{4}{c}{Ordered per faculty} \\
\cmidrule(r{1.2em}){2-4}\cmidrule(l){5-8}
\# & Department & Score & Per-fac.\ & Department & Per fac.\ & Faculty & \csr{} \\
\midrule
1  & CMU          & 20.6 & 7  & Stanford     & 11.66 & 61  & 10 \\
2  & UIUC         & 16.7 & 10 & UC Berkeley  & 8.32  & 85  & 7  \\
3  & UCSD         & 14.8 & 14 & UCLA         & 7.55  & 50  & 20 \\
4  & Georgia Tech & 13.7 & 30 & MIT          & 7.50  & 104 & 5  \\
5  & MIT          & 11.3 & 4  & UT Austin    & 7.28  & 60  & 14 \\
6  & UW           & 10.7 & 6  & UW           & 7.04  & 80  & 6  \\
7  & UC Berkeley  & 10.6 & 2  & CMU          & 6.59  & 183 & 1  \\
8  & UMD          & 10.3 & 12 & Penn         & 6.31  & 70  & 18 \\
9  & UMich        & 10.1 & 23 & Harvard      & 6.22  & 38  & 39 \\
10 & Stanford     & 9.6  & 1  & UIUC         & 6.16  & 127 & 2  \\
11 & Northeastern & 9.4  & 39 & Princeton    & 5.90  & 67  & 17 \\
12 & Purdue       & 8.5  & 29 & UMD          & 5.88  & 98  & 8  \\
13 & Cornell      & 8.2  & 19 & Columbia     & 5.83  & 58  & 19 \\
14 & NYU          & 7.7  & 17 & UCSD         & 5.62  & 130 & 3  \\
15 & UT Austin    & 7.7  & 5  & UNC          & 5.58  & 46  & 43 \\
\bottomrule
\end{tabular}
\end{table}

The two orderings agree far less than the venue sets did. Kendall's $\tau$
between them is 0.716, against 0.95 or higher for every pair of venue sets in
Table~\ref{tab:venue}; the median department moves 13 positions and one moves
71. Only 6 of the 10 departments in the \csr{} top ten remain in the top ten
when measured per head. Figure~\ref{fig:percapita} plots the disagreement
directly.

\begin{figure}[t]
\centering
\begin{tikzpicture}
\begin{axis}[cspicks, width=0.62\textwidth, height=6.4cm,
  xlabel={\csr{} rank}, ylabel={Per-faculty rank},
  xmin=0, xmax=170, ymin=0, ymax=170]
\addplot[gray!55, dashed, thick, forget plot] coordinates {(0,0) (170,170)};
\addplot[only marks, mark=*, mark size=1.1pt, blue!70!black, opacity=0.55]
  table [x=csrank, y=pcrank] {data/percapita-scatter.dat};
\end{axis}
\end{tikzpicture}
\caption{Departmental rank under the \csr{} score against rank by adjusted
output per publishing faculty member, for the 166 U.S.\ departments with at
least five publishing faculty. Points below the diagonal do better per head
than in aggregate. $\tau = 0.716$.}
\label{fig:percapita}
\end{figure}

The direction of the change is systematic. CMU leads the \csr{} ordering with
183 publishing faculty and places seventh per head; Stanford places tenth
overall and first per head with 61. Georgia Tech is the sharpest case inside
the top five: fourth by \csr{} score and thirtieth per head. Northeastern,
eleventh overall, is thirty-ninth. Running the other way, Harvard is 39th
overall and 9th per head on 38 faculty, and UNC is 43rd overall and 15th per
head on 46.

Table~\ref{tab:percapitagains} lists the departments that gain most from the
change of question. They are small and selective, and the volume-weighted
ordering buries them.

\begin{table}[t]
\centering
\caption{Departments gaining the most positions when ordered per publishing
faculty member rather than by \csr{} score.}
\label{tab:percapitagains}
\begin{tabular}{lrrr}
\toprule
Department & \csr{} rank & Per-faculty rank & Publishing faculty \\
\midrule
TTI Chicago    & 87  & 16 & 15 \\
Caltech        & 80  & 20 & 27 \\
UC Merced      & 73  & 18 & 18 \\
New Hampshire  & 110 & 60 & 8  \\
New Mexico     & 127 & 79 & 8  \\
\bottomrule
\end{tabular}
\end{table}

Neither ordering is more correct. A large department genuinely produces more
research, which is what \csr{} sets out to measure, and an applicant choosing
by prestige or by breadth of available advisors is right to weigh size. But an
applicant asking how productive the people around them will be is asking the
per-head question, and the volume-weighted ordering does not answer it. Our
point is that the two questions have materially different answers --- 13
positions of median disagreement --- and that a single ordered list conceals
which one it is reporting.

\subsection{Many departments rest on very few people}
\label{sec:fragilityresults}

Table~\ref{tab:fragility} reports how many departures would move a department
out of its own rank band, for the 51 departments ranked in the top 50 with at
least five publishing faculty.

\begin{table}[t]
\centering
\caption{Departures required to leave a rank band, by band. Three departments do
not leave their band within the 15-removal search and are excluded from the
medians.}
\label{tab:fragility}
\begin{tabular}{lrr}
\toprule
Band & Departments & Median departures \\
\midrule
Top 10  & 8  & 3.5 \\
Top 25  & 15 & 5 \\
Top 50  & 25 & 2 \\
\midrule
All resolved & 48 & 3 \\
\bottomrule
\end{tabular}
\end{table}

Twelve of the 48 departments leave their band on a \emph{single} departure, and
22 --- almost half --- on two or fewer. The median is three.
Table~\ref{tab:fragileexamples} gives named cases.

\begin{table}[t]
\centering
\caption{Departures required to leave a rank band, selected departments.
``$>$15'' means the department does not leave its band within the search
horizon. Top-3 share is the fraction of departmental adjusted output produced
by its three most prolific members; the individuals are not identified
(\S\ref{sec:ethics}).}
\label{tab:fragileexamples}
\begin{tabular}{lrrrr}
\toprule
Department & Rank & Band & Publishing faculty & Departures \\
\midrule
CMU          & 1  & 10 & 183 & $>$15 \\
MIT          & 5  & 10 & 104 & 5 \\
Northeastern & 11 & 25 & 115 & $>$15 \\
GMU          & 32 & 50 & 54  & 3 \\
Rice         & 43 & 50 & 39  & 1 \\
\bottomrule
\end{tabular}
\end{table}

Rice, ranked 43rd with 39 publishing faculty, falls out of the top 50 when one
person leaves. GMU, 32nd with 54 faculty, needs three. At the other end,
CMU and Northeastern never leave their bands within the 15-removal
search, and MIT needs five departures to leave the top ten.

Fragility tracks concentration, as it should: departments that leave their band
on one departure have a median top-3 share of 22.1\%, against 15.4\% for those
needing four or more. The band medians are not monotonic --- the top-50 band has
the lowest median at 2 --- because a department sitting near the bottom edge of
a band is close to leaving it regardless of how its output is distributed.

Read alongside \S\ref{sec:stabilityresults}, this sharpens the picture. A rank
is contingent on methodology, moving a median of 17 positions across defensible
settings; it is also contingent on a handful of individuals continuing to work
where they currently do. For an applicant, a department whose standing rests on
one or two people is a genuinely different proposition from one of the same rank
whose output is spread across dozens, and the rank alone does not distinguish
them.

\subsection{Individual output is more concentrated than departmental output}
\label{sec:researchers}

Departments are aggregates of people, and the aggregate hides a distribution
far more skewed than the departmental one. Table~\ref{tab:researcherdist}
summarizes the 5{,}987 U.S.-affiliated researchers with any credited output in
the 10-year window.

\begin{table}[t]
\centering
\caption{The distribution of adjusted output and breadth across 5{,}987
U.S.-affiliated researchers with credited output in a 10-year window.
Specialists take at least 60\% of their output in one area; generalists are
active in at least four.}
\label{tab:researcherdist}
\begin{tabular}{lr}
\toprule
Quantity & Value \\
\midrule
Researchers with credited output    & 5{,}987 \\
Median adjusted count               & 2.19 \\
Median breadth (areas)              & 2 \\
\midrule
Share of output from the top 1\%    & 8.6\% \\
Share of output from the top 5\%    & 25.9\% \\
Share of output from the top 10\%   & 40.0\% \\
\midrule
Active in exactly one area          & 28.8\% \\
Specialists ($\ge$60\% in one area) & 69.4\% \\
Generalists ($\ge$4 areas)          & 29.3\% \\
\bottomrule
\end{tabular}
\end{table}

The top decile of researchers produces two fifths of all U.S.\ output, and the
median researcher contributes an adjusted count of 2.19 over ten years --- the
equivalent of roughly one four-author paper every two years. This is the
individual-level counterpart of \S\ref{sec:concentration}, and it is the
mechanism behind it: departmental concentration is high because the underlying
individual distribution is heavy-tailed, and a department is simply a small
sample from it. It also explains why fragility is so often resolved in one or
two departures.

Breadth is the other surprise. The median researcher is active in two areas and
nearly seven in ten take the majority of their output in a single one, which
sits awkwardly beside a departmental score (Eq.~\ref{eq:score}) that rewards
breadth across 27 of them. Departmental breadth is assembled from specialists,
not produced by generalists.

Table~\ref{tab:topresearchers} lists the most prolific individuals.

\begin{table}[t]
\centering
\small
\caption{The fifteen U.S.-affiliated researchers with the highest adjusted
count over a 10-year window. ``Share'' is the fraction of adjusted output in
the researcher's largest area. All quantities are derived from publication
counts \csr{} already publishes; see \S\ref{sec:ethics}.}
\label{tab:topresearchers}
\begin{tabular}{rllrrrlr}
\toprule
\# & Researcher & Department & Adj. & Papers & Areas & Primary area & Share \\
\midrule
1  & Sergey Levine     & UC Berkeley   & 78.9 & 370 & 6 & Machine Learning & 71\% \\
2  & Mohit Bansal      & UNC           & 57.6 & 216 & 5 & NLP              & 61\% \\
3  & Pieter Abbeel     & UC Berkeley   & 53.8 & 246 & 8 & Machine Learning & 66\% \\
4  & Heng Huang        & UMD           & 51.3 & 186 & 8 & Machine Learning & 42\% \\
5  & Stefano Ermon     & Stanford      & 47.5 & 201 & 3 & Machine Learning & 80\% \\
6  & Graham Neubig     & CMU           & 46.5 & 201 & 7 & NLP              & 70\% \\
7  & David P. Woodruff & CMU           & 46.0 & 146 & 9 & Machine Learning & 57\% \\
8  & Ming-Hsuan Yang   & UC Merced     & 45.5 & 254 & 6 & Computer Vision  & 76\% \\
9  & Dinesh Manocha    & UMD           & 44.3 & 210 & 9 & Robotics         & 47\% \\
10 & Daniela Rus       & MIT           & 42.5 & 201 & 7 & Robotics         & 74\% \\
11 & Quanquan Gu       & UCLA          & 40.3 & 155 & 6 & Machine Learning & 91\% \\
12 & Trevor Darrell    & UC Berkeley   & 39.1 & 210 & 6 & Computer Vision  & 50\% \\
13 & Alan L. Yuille    & Johns Hopkins & 36.7 & 193 & 5 & Computer Vision  & 70\% \\
14 & J. Zico Kolter    & CMU           & 35.8 & 134 & 4 & Machine Learning & 87\% \\
15 & Kristen Grauman   & UT Austin     & 34.6 & 107 & 4 & Computer Vision  & 77\% \\
\bottomrule
\end{tabular}
\end{table}

Two observations. The head of the distribution is entirely AI-adjacent: all
fifteen take their primary area in machine learning (seven), computer vision
(four), NLP (two), or robotics (two). This is a fact about where publication
volume is (\S\ref{sec:fields}) and not a statement about research importance.
And adjusted count is not paper count: Ming-Hsuan Yang's 254 papers exceed
those of every researcher ranked above him but one, yet he places eighth on
adjusted count, because fractional credit divides a paper among its authors and
the two orderings diverge for large-collaboration work. A tool that reported
raw paper counts would produce a visibly different list.

The most prolific researchers are also, mostly, specialists. The widest
researchers are a different set, shown in Table~\ref{tab:broadest}.

\begin{table}[t]
\centering
\caption{The widest-ranging researchers among the 300 most prolific, by number
of top-level areas with credited output, 10-year window.}
\label{tab:broadest}
\begin{tabular}{llrrl}
\toprule
Researcher & Department & Areas & Adj. & Primary area \\
\midrule
Ion Stoica     & UC Berkeley  & 12 & 20.1 & Machine Learning \\
Yanzhi Wang    & Northeastern & 12 & 15.9 & Design Automation \\
Cong Liu       & UC Riverside & 12 & 14.1 & Embedded \& Real-Time \\
Isil Dillig    & UT Austin    & 12 & 13.6 & Programming Languages \\
Dawn Song      & UC Berkeley  & 11 & 25.3 & Machine Learning \\
Xiangyu Zhang  & Purdue       & 10 & 20.3 & Software Engineering \\
Yiran Chen     & Duke         & 10 & 18.5 & Design Automation \\
Osbert Bastani & Penn         & 10 & 15.3 & Machine Learning \\
\bottomrule
\end{tabular}
\end{table}

None of these appear in Table~\ref{tab:topresearchers}, and none is close to
the top on volume --- twelve areas at an adjusted count of 20.1 is a quarter of
the leader's output spread twice as wide. Under Eq.~\ref{eq:score} a department
values these two profiles very differently: because the geometric mean
multiplies $(c+1)$ over \emph{all} areas, a researcher who supplies the only
credit in an otherwise empty area moves the departmental score more than a
larger producer in a crowded one. This is the individual-level statement of the
same effect that decides the departmental comparison in
\S\ref{sec:comparison}.

Table~\ref{tab:arealeaders} gives the most prolific researcher in each of the
27 subfields, which is the query an applicant with a settled subfield actually
wants and which no departmental ordering answers.

\begin{table}[p]
\centering
\small
\caption{The most prolific U.S.-affiliated researcher in each top-level
subfield, 10-year window, with their adjusted count \emph{within that subfield}
and their overall position in the researcher ordering. Subfields are ordered by
region-wide output.}
\label{tab:arealeaders}
\begin{tabular}{llrlr}
\toprule
Subfield & Researcher & In-area adj. & Department & Overall \# \\
\midrule
Machine Learning      & Sergey Levine          & 55.8 & UC Berkeley   & 1 \\
Computer Vision       & Ming-Hsuan Yang        & 34.4 & UC Merced     & 8 \\
NLP                   & Mohit Bansal           & 34.9 & UNC           & 2 \\
AI                    & Heng Huang             & 18.0 & UMD           & 4 \\
Robotics              & Daniela Rus            & 31.6 & MIT           & 10 \\
HCI                   & Chris Harrison         & 18.5 & CMU           & 111 \\
Algorithms            & Thatchaphol Saranurak  & 18.3 & UMich         & 117 \\
Security              & Zhiqiang Lin           & 12.6 & Ohio State    & 165 \\
Databases             & Immanuel Trummer       & 15.9 & Cornell       & 160 \\
Architecture          & Moinuddin K. Qureshi   & 11.8 & Georgia Tech  & 236 \\
Info Retrieval        & Hamed Zamani           & 12.2 & UMass Amherst & 242 \\
Design Automation     & David Z. Pan           & 10.8 & UT Austin     & 300 \\
Software Engineering  & Hridesh Rajan          & 8.2  & Tulane        & 672 \\
Cryptography          & Stefano Tessaro        & 16.0 & UW            & 127 \\
Graphics              & Ravi Ramamoorthi       & 10.3 & UCSD          & 106 \\
Visualization         & Kwan-Liu Ma            & 7.3  & UC Davis      & 580 \\
Programming Languages & Isil Dillig            & 5.9  & UT Austin     & 259 \\
Mobile Computing      & Xinyu Zhang            & 11.7 & UCSD          & 107 \\
Networks              & Mohammad Alizadeh      & 6.8  & MIT           & 354 \\
HPC                   & Devesh Tiwari          & 6.3  & Northeastern  & 318 \\
CS Education          & Craig B. Zilles        & 7.1  & UIUC          & 808 \\
Measurement           & Ziv Scully             & 4.1  & Cornell       & 1376 \\
Embedded \& Real-Time & James H. Anderson      & 9.6  & UNC           & 515 \\
Econ \& Computation   & S. Matthew Weinberg    & 7.8  & Princeton     & 196 \\
Operating Systems     & Nickolai Zeldovich     & 4.8  & MIT           & 1191 \\
Logic \& Verification & Kuldeep S. Meel        & 4.9  & Georgia Tech  & 70 \\
Comp.\ Biology        & Carl Kingsford         & 5.9  & CMU           & 837 \\
\bottomrule
\end{tabular}
\end{table}

The overall-rank column is the point of the table. The leader of a subfield is
frequently nowhere near the top of the aggregate researcher ordering: the most
prolific U.S.\ researcher in measurement sits 1376th overall, in operating
systems 1191st, in computational biology 837th, and in CS education 808th.
Only 6 of the 27 subfield leaders appear in the overall top 100, and the median
subfield leader ranks 196th. This is not a statement
about those researchers; it is a statement about the aggregate ordering, which
is dominated by the volume differences between subfields
(\S\ref{sec:fields}) rather than by standing within one. An applicant who
searches a subfield sees the first column; an applicant who reads a global list
of researchers sees the last, and the two disagree almost completely. The same
caution applies to the ordering in Table~\ref{tab:topresearchers}.

\subsection{Subfields differ by two orders of magnitude, and the gap is widening}
\label{sec:fields}

Table~\ref{tab:fields} reports every top-level subfield at region-wide scale.

\begin{table}[p]
\centering
\small
\caption{U.S.\ output by subfield over a 10-year window, with growth against
the equal-length preceding period. ``Depts'' and ``Researchers'' count distinct
entities with any credit in the subfield; ``Top-5'' is the share of the
subfield's national output held by its five largest departments.}
\label{tab:fields}
\setlength{\tabcolsep}{4.5pt}
\begin{tabular}{lrrrrrlr}
\toprule
Subfield & Adjusted & Growth & Depts & Researchers & Per res.\ & Leader & Top-5 \\
\midrule
Machine Learning      & 5382.7 & $+444$\% & 165 & 2422 & 2.22 & CMU          & 24.2\% \\
Computer Vision       & 2076.2 & $+83$\%  & 156 & 1185 & 1.75 & CMU          & 20.4\% \\
NLP                   & 1919.9 & $+189$\% & 154 & 1175 & 1.63 & CMU          & 22.1\% \\
AI                    & 1806.4 & $+146$\% & 172 & 1861 & 0.97 & CMU          & 14.5\% \\
Robotics              & 1707.6 & $+52$\%  & 144 & 906  & 1.89 & CMU          & 24.0\% \\
HCI                   & 1476.3 & $+146$\% & 146 & 1202 & 1.23 & CMU          & 28.5\% \\
Algorithms            & 1241.7 & $+20$\%  & 91  & 480  & 2.59 & CMU          & 30.2\% \\
Security              & 1137.9 & $+164$\% & 149 & 1039 & 1.10 & Georgia Tech & 19.3\% \\
Databases             & 573.2  & $+49$\%  & 125 & 570  & 1.01 & MIT          & 20.6\% \\
Architecture          & 565.2  & $+82$\%  & 101 & 608  & 0.93 & UIUC         & 26.7\% \\
Info Retrieval        & 438.7  & $+75$\%  & 140 & 729  & 0.60 & UIUC         & 22.8\% \\
Design Automation     & 407.5  & $+10$\%  & 118 & 412  & 0.99 & UCSD         & 25.4\% \\
Software Engineering  & 397.0  & $+60$\%  & 118 & 417  & 0.95 & CMU          & 25.5\% \\
Cryptography          & 336.0  & $+72$\%  & 64  & 160  & 2.10 & UW           & 32.6\% \\
Graphics              & 315.6  & $+10$\%  & 73  & 229  & 1.38 & CMU          & 36.5\% \\
Visualization         & 310.4  & $+76$\%  & 110 & 337  & 0.92 & Utah         & 23.5\% \\
Programming Languages & 309.1  & $+47$\%  & 91  & 363  & 0.85 & CMU          & 25.6\% \\
Mobile Computing      & 301.2  & $+51$\%  & 105 & 413  & 0.73 & UCSD         & 22.2\% \\
Networks              & 290.4  & $+29$\%  & 91  & 405  & 0.72 & Princeton    & 32.3\% \\
HPC                   & 282.0  & $-4$\%   & 123 & 444  & 0.63 & Ohio State   & 18.6\% \\
CS Education          & 264.4  & $+251$\% & 107 & 336  & 0.79 & UCSD         & 28.7\% \\
Measurement           & 220.5  & $+10$\%  & 91  & 349  & 0.63 & Northeastern & 25.5\% \\
Embedded \& Real-Time & 160.4  & $-27$\%  & 76  & 164  & 0.98 & WUSTL        & 38.1\% \\
Econ \& Computation   & 154.4  & $-27$\%  & 55  & 142  & 1.09 & Cornell      & 30.6\% \\
Operating Systems     & 147.9  & $+85$\%  & 65  & 274  & 0.54 & UC Berkeley  & 33.5\% \\
Logic \& Verification & 119.6  & $+18$\%  & 67  & 159  & 0.75 & Georgia Tech & 28.1\% \\
Comp.\ Biology        & 105.0  & $-31$\%  & 78  & 166  & 0.63 & CMU          & 33.9\% \\
\bottomrule
\end{tabular}
\end{table}

\begin{figure}[t]
\centering
\begin{tikzpicture}
\begin{axis}[cspicks, width=0.72\textwidth, height=11cm,
  xbar, bar width=6pt,
  xlabel={Adjusted count}, ylabel={},
  xmin=0, ymin=0.3, ymax=27.7, y dir=reverse,
  ytick=data, yticklabels from table={data/field-totals.dat}{label},
  y tick label style={font=\footnotesize},
  xmajorgrids, ymajorgrids=false,
  nodes near coords, nodes near coords style={font=\tiny, /pgf/number format/precision=0, /pgf/number format/fixed}]
\addplot[fill=blue!55!black, draw=none] table [x=total, y=idx] {data/field-totals.dat};
\end{axis}
\end{tikzpicture}
\caption{U.S.\ adjusted publication count by top-level subfield over a 10-year
window. Machine learning alone accounts for 24.0\% of the total; the four
AI-adjacent subfields together account for 49.8\%.}
\label{fig:fields}
\end{figure}

The distribution is extraordinarily skewed (Figure~\ref{fig:fields}). Machine
learning alone accounts for 24.0\% of all U.S.\ adjusted output. Together with
computer vision, NLP, and AI it accounts for 49.8\% --- half the measured
field. At the other end, computational biology, logic and verification, and
operating systems each account for well under 1\%. The ratio between the
largest and smallest subfield is 51:1.

The growth column shows this is recent and accelerating. Machine learning grew
444\% against the preceding decade, NLP 189\%, security 164\%. Three subfields
contracted outright: computational biology by 31\%, embedded and real-time
systems and economics-and-computation by 27\% each. CS education grew 251\%
from a small base, spreading to 107 departments.

This has a direct consequence for how Eq.~\ref{eq:score} behaves that is easy
to miss. The geometric mean gives every one of the 27 areas equal weight in the
exponent, but the areas are not equally sized, so a unit of adjusted output is
worth far more to a department's score in operating systems --- where the
national total is 147.9 and $\ln(c+1)$ is steeply increasing --- than in
machine learning, where a department must add enormous volume to move the same
factor. A department that hires one systems researcher may gain more score than
one that hires three machine-learning researchers. This is arguably the correct
incentive for a breadth-rewarding measure, but it is nowhere visible in the
output, and it is why per-area credit and not just totals must be shown.

Subfield leadership also turns over. Table~\ref{tab:fieldlead} lists the
subfields whose leading department changed against the preceding period.

\begin{table}[t]
\centering
\caption{Subfields whose leading U.S.\ department changed between the preceding
decade and the current one.}
\label{tab:fieldlead}
\begin{tabular}{llll r}
\toprule
Subfield & Former leader & Current leader & & Growth \\
\midrule
NLP          & Stanford & CMU          & & $+189$\% \\
HCI          & UW       & CMU          & & $+146$\% \\
Algorithms   & MIT      & CMU          & & $+20$\% \\
Databases    & CMU      & MIT          & & $+49$\% \\
Architecture & UMich    & UIUC         & & $+82$\% \\
Cryptography & UCLA     & UW           & & $+72$\% \\
\bottomrule
\end{tabular}
\end{table}

CMU now leads 11 of the 27 subfields, having taken NLP, HCI, and algorithms
from three different departments while losing databases to MIT. This is the
concrete content behind its invariant first place in
Table~\ref{tab:examples}: it is not first because it is far ahead anywhere in
particular, but because it is at or near the front almost everywhere, which is
exactly what Eq.~\ref{eq:score} rewards.

Finally, the interface compares two subfields directly.
Table~\ref{tab:fieldcompare} gives machine learning against programming
languages, chosen as an extreme pair.

\begin{table}[t]
\centering
\caption{Two subfields head to head, 10-year window against the preceding
decade. ``New researchers'' are those with credit in the subfield now and none
in the preceding period.}
\label{tab:fieldcompare}
\begin{tabular}{lrr}
\toprule
& Machine Learning & Programming Languages \\
\midrule
Adjusted count            & 5382.7 & 309.1 \\
Preceding decade          & 989.1  & 210.0 \\
Growth                    & $+444$\% & $+47$\% \\
Departments active        & 165    & 91  \\
Researchers active        & 2422   & 363 \\
New researchers           & 1857   & 206 \\
\midrule
Leading department        & CMU (315.1) & CMU (24.1) \\
Second                    & MIT (267.1) & MIT (15.6) \\
Third                     & UC Berkeley (264.3) & Penn (14.3) \\
\midrule
Researchers active in both & \multicolumn{2}{c}{113} \\
\bottomrule
\end{tabular}
\end{table}

The two subfields are separated by a factor of 17 in output and a factor of 6.7
in participation, and 1{,}857 researchers entered machine learning in the last
decade --- more than five times the entire current population of programming
languages. Yet under Eq.~\ref{eq:score} they contribute symmetric factors to a
department's score. An applicant told that a department is ``ranked 12th''
learns nothing about which side of this divide its strength lies on, which is
why the per-subfield view (\S\ref{sec:attribution}) rather than the rank is the
thing we put in front of them.

\subsection{Head-to-head comparison isolates what a rank gap is made of}
\label{sec:comparison}

Table~\ref{tab:compareuiuc} compares UIUC and MIT, three positions apart in the
current ranking.

\begin{table}[t]
\centering
\caption{UIUC against MIT, 10-year window. ``Areas led'' counts top-level areas
where the department has strictly more adjusted credit.}
\label{tab:compareuiuc}
\begin{tabular}{lrr}
\toprule
& UIUC & MIT \\
\midrule
Overall rank        & 2      & 5 \\
Papers              & 3488   & 3294 \\
Adjusted count      & 782.1  & 779.7 \\
Publishing faculty  & 127    & 104 \\
Areas led           & 17     & 10 \\
\midrule
\multicolumn{3}{l}{\emph{Largest per-area margins}} \\
Machine Learning    & 195.8 & \textbf{267.1} \\
Algorithms          & 37.8  & \textbf{86.3} \\
Robotics            & 37.0  & \textbf{85.5} \\
NLP                 & \textbf{72.1} & 34.3 \\
HCI                 & \textbf{42.6} & 22.7 \\
\midrule
\multicolumn{3}{l}{\emph{Areas that produced the rank gap} (Eq.~\ref{eq:gap})} \\
Info Retrieval      & \textbf{33.1} & 2.3 \\
CS Education        & \textbf{14.5} & 0.6 \\
Security            & \textbf{48.8} & 8.8 \\
HPC                 & \textbf{6.4}  & 0.5 \\
Embedded \& Real-Time & \textbf{3.5} & 0.0 \\
\bottomrule
\end{tabular}
\end{table}

This is the clearest single demonstration of what Eq.~\ref{eq:score} actually
measures. The two departments' total adjusted output differs by 2.4 --- 0.3\%
--- and MIT wins the three largest per-area margins outright, taking machine
learning by 71.3, algorithms by 48.5, and robotics by 48.5. Yet UIUC ranks
three positions higher.

The decomposition explains it. Every area in the bottom block is one where MIT
has almost no credited output: 2.3 in information retrieval, 0.6 in CS
education, 0.0 in embedded and real-time systems. Because the score multiplies
$(c+1)$ across all 27 areas, the move from 0.0 to 3.5 contributes more to the
product than the move from 195.8 to 267.1 --- $\ln(4.5)-\ln(1.0) = 1.50$ against
$\ln(268.1)-\ln(196.8) = 0.31$. The rank gap is made almost entirely of areas
where neither department is nationally dominant and one of them is simply
absent.

Neither the raw margins nor the rank alone would tell a reader this. A
prospective machine learning student reading ``UIUC is ranked above MIT'' is
being told something true about a breadth-weighted aggregate and something
false about the thing they care about, where MIT leads by 36\%. This is the
central case for reporting per-area credit alongside the ordinal.

Table~\ref{tab:comparegmu} gives a second pair further down the ranking, where
the same mechanism operates.

\begin{table}[t]
\centering
\caption{GMU against Rice, 10-year window.}
\label{tab:comparegmu}
\begin{tabular}{lrr}
\toprule
& GMU & Rice \\
\midrule
Overall rank        & 32    & 43 \\
Papers              & 776   & 678 \\
Adjusted count      & 178.7 & 158.3 \\
Publishing faculty  & 54    & 39 \\
Areas led           & 15    & 12 \\
\midrule
\multicolumn{3}{l}{\emph{Largest per-area margins}} \\
Machine Learning    & 14.5 & \textbf{37.8} \\
AI                  & 9.3  & \textbf{23.8} \\
Computer Vision     & 5.5  & \textbf{17.3} \\
HCI                 & \textbf{22.1} & 2.1 \\
Security            & \textbf{22.6} & 3.6 \\
NLP                 & \textbf{22.2} & 4.4 \\
\midrule
\multicolumn{3}{l}{\emph{Areas that produced the rank gap} (Eq.~\ref{eq:gap})} \\
HCI                 & \textbf{22.1} & 2.1 \\
Visualization       & \textbf{5.3}  & 0.0 \\
Security            & \textbf{22.6} & 3.6 \\
Software Engineering & \textbf{7.7} & 0.8 \\
Graphics            & \textbf{5.5}  & 0.3 \\
\bottomrule
\end{tabular}
\end{table}

Rice is the stronger machine learning and AI department by a wide margin ---
2.6$\times$ in machine learning --- and ranks eleven positions lower. GMU's rank
advantage is built from HCI, security, visualization, software engineering, and
graphics, five areas in which Rice is near zero. Recall from
\S\ref{sec:fragilityresults} that Rice leaves the top 50 on a single departure
while GMU needs three: the two comparisons together say that these departments
are not close substitutes at any level of detail, despite sitting eleven
positions apart in an ordering that reduces both to one number.

The same view works for individuals. Table~\ref{tab:compareresearcher} compares
the two most prolific U.S.\ researchers in the window.

\begin{table}[t]
\centering
\caption{Two researchers head to head, 10-year window.}
\label{tab:compareresearcher}
\begin{tabular}{lrr}
\toprule
& Sergey Levine & Mohit Bansal \\
\midrule
Department       & UC Berkeley & UNC \\
Papers           & 370   & 216 \\
Adjusted count   & 78.9  & 57.6 \\
Active areas     & 6     & 5 \\
Areas led        & 3     & 3 \\
\midrule
Machine Learning & \textbf{55.8} & 11.2 \\
NLP              & 0.8   & \textbf{34.9} \\
Robotics         & \textbf{19.8} & 0.5 \\
Computer Vision  & 1.5   & \textbf{5.8} \\
AI               & 0.3   & \textbf{5.2} \\
\bottomrule
\end{tabular}
\end{table}

The scoreboard reports a clear leader on volume and a tie on areas led, which
is the honest summary: these are researchers in adjacent but different
subfields, and the aggregate that puts one above the other is dominated by a
single area in which they barely overlap. The interface says so rather than
declaring a winner, and this is the reason the comparison view reports areas
led alongside the total instead of collapsing to one verdict.

\subsection{Funding orders departments differently from publications}
\label{sec:nsf}

The NSF snapshot provides a second, independently produced signal over the same
roster. Table~\ref{tab:nsfschools} orders U.S.\ departments by attributed NSF
award value over the same 10-year window and reports their publication rank
alongside.

\begin{table}[t]
\centering
\small
\caption{The fifteen U.S.\ departments with the largest attributed NSF award
value, awards dated 2016--2026. ``Attributed'' splits each award equally among
its listed investigators (\S\ref{sec:nsfmethod}). ``Pub.\ rank'' is the \csr{}
rank over the same window.}
\label{tab:nsfschools}
\begin{tabular}{rlrrrrr}
\toprule
\# & Department & Attributed & Awards & Funded faculty & Pub.\ rank & $\Delta$ \\
\midrule
1  & CMU          & \$142.7M & 293 & 124 & 1  & 0 \\
2  & Georgia Tech & \$119.9M & 278 & 103 & 4  & $+2$ \\
3  & UCSD         & \$119.1M & 235 & 97  & 3  & 0 \\
4  & UIUC         & \$108.3M & 236 & 96  & 2  & $-2$ \\
5  & Northeastern & \$107.0M & 216 & 75  & 11 & $+6$ \\
6  & UT Austin    & \$91.4M  & 109 & 40  & 14 & $+8$ \\
7  & UW           & \$82.3M  & 171 & 58  & 6  & $-1$ \\
8  & MIT          & \$78.5M  & 173 & 69  & 5  & $-3$ \\
9  & UW--Madison  & \$75.8M  & 157 & 58  & 14 & $+5$ \\
10 & UMD          & \$75.5M  & 199 & 78  & 8  & $-2$ \\
11 & UMich        & \$73.2M  & 179 & 69  & 9  & $-2$ \\
12 & Northwestern & \$72.5M  & 153 & 54  & 25 & $+13$ \\
13 & NYU          & \$70.8M  & 166 & 55  & 14 & $+1$ \\
14 & Rice         & \$63.0M  & 83  & 28  & 43 & $+29$ \\
15 & Virginia     & \$62.3M  & 151 & 47  & 28 & $+13$ \\
\bottomrule
\end{tabular}
\end{table}

The head of the two orderings agrees closely --- the same four departments lead
both --- but agreement decays quickly. Across the 182 matched departments,
Kendall's $\tau$ between funding rank and publication rank is 0.778, and the
median department sits 12 positions apart in the two orderings. That is worse
agreement than any pair of venue sets in Table~\ref{tab:venue} and comparable
to the per-faculty disagreement of \S\ref{sec:percapitaresults}.
Figure~\ref{fig:funding} plots it.

\begin{figure}[t]
\centering
\begin{tikzpicture}
\begin{axis}[cspicks, width=0.62\textwidth, height=6.4cm,
  xlabel={Publication rank}, ylabel={NSF funding rank},
  xmin=0, xmax=185, ymin=0, ymax=185]
\addplot[gray!55, dashed, thick, forget plot] coordinates {(0,0) (185,185)};
\addplot[only marks, mark=*, mark size=1.1pt, blue!70!black, opacity=0.55]
  table [x=pubrank, y=fundrank] {data/funding-scatter.dat};
\end{axis}
\end{tikzpicture}
\caption{Publication rank against attributed-NSF-funding rank for the 182
matched U.S.\ departments, 10-year window. $\tau = 0.778$, median displacement
12 positions.}
\label{fig:funding}
\end{figure}

The divergent cases are informative in both directions
(Table~\ref{tab:nsfgap}). Rice is 43rd on publications and 14th on funding.
Texas A\&M is 34th on publications and 167th on funding --- a 133-position gap,
by far the largest in the data, and one worth treating as a question about the
matching rather than a finding about the department (\S\ref{sec:limitations}).

\begin{table}[t]
\centering
\caption{Departments where the funding and publication orderings disagree most,
in each direction. Gap is publication rank minus funding rank.}
\label{tab:nsfgap}
\begin{tabular}{lrrr}
\toprule
Department & Funding rank & Pub.\ rank & Gap \\
\midrule
\multicolumn{4}{l}{\emph{Funded well above publication standing}} \\
Rhode Island      & 104 & 159 & $+55$ \\
New Mexico State  & 86  & 139 & $+53$ \\
Missouri          & 88  & 139 & $+51$ \\
UNLV              & 108 & 159 & $+51$ \\
Hawaii at Manoa   & 103 & 139 & $+36$ \\
\midrule
\multicolumn{4}{l}{\emph{Publishing well above funding standing}} \\
Texas A\&M        & 167 & 34  & $-133$ \\
Brigham Young     & 154 & 110 & $-44$ \\
Arizona           & 122 & 82  & $-40$ \\
UC Merced         & 110 & 73  & $-37$ \\
Old Dominion      & 176 & 139 & $-37$ \\
\bottomrule
\end{tabular}
\end{table}

Individual attribution is far more concentrated than departmental attribution,
and far more sensitive to the equal-split assumption
(Table~\ref{tab:nsffaculty}).

\begin{table}[t]
\centering
\small
\caption{The twelve investigators with the largest attributed NSF share,
2016--2026. ``Project value'' is the full value of the awards they participate
in, which for large collaborative awards is many times the attributed share.}
\label{tab:nsffaculty}
\begin{tabular}{rllrrr}
\toprule
\# & Investigator & Department & Attributed & Project value & Awards \\
\midrule
1  & Adam R. Klivans           & UT Austin        & \$41.2M & \$42.3M  & 5  \\
2  & Richard G. Baraniuk       & Rice             & \$21.3M & \$109.3M & 8  \\
3  & Sidney K. D'Mello         & Colorado Boulder & \$11.0M & \$47.1M  & 10 \\
4  & Ashok K. Goel             & Georgia Tech     & \$10.8M & \$20.8M  & 4  \\
5  & John Preskill             & Caltech          & \$10.7M & \$25.3M  & 2  \\
6  & James C. Lester           & NC State         & \$10.4M & \$38.1M  & 11 \\
7  & Ellie Pavlick             & Brown            & \$10.3M & \$21.0M  & 2  \\
8  & Suresh Venkatasubramanian & Brown            & \$10.0M & \$20.0M  & 1  \\
9  & Tamara Sumner             & Colorado Boulder & \$9.4M  & \$47.9M  & 7  \\
10 & David R. Choffnes         & Northeastern     & \$8.9M  & \$22.2M  & 12 \\
11 & Tiffany Barnes            & NC State         & \$7.4M  & \$23.3M  & 18 \\
12 & Yiran Chen                & Duke             & \$7.2M  & \$30.2M  & 20 \\
\bottomrule
\end{tabular}
\end{table}

The gap between the two money columns is the equal-split assumption made
visible. The top investigator's five awards are almost entirely his own on
paper (\$41.2M attributed against \$42.3M of project value), while the second's
eight awards carry \$109.3M of project value of which the equal split assigns
him \$21.3M. Neither number is the truth about how much money either person
directs, and the interface shows both precisely so that neither can be read as
if it were. Note also that this list has almost no overlap with the publication
leaders of Table~\ref{tab:topresearchers}: of the twelve, only Yiran Chen
appears anywhere in the breadth or output tables above. Funding and publication
volume are close to orthogonal at the individual level.

The award population itself is worth characterizing. Of the 10{,}226 awards in
the window, 81.6\% are from the CISE directorate, with education (486),
engineering (463), and mathematical and physical sciences (284) making up most
of the remainder --- so this is not a CISE-only view, and a department strong in
interdisciplinary work will show funding that its publication record does not
predict. The largest programs by award count are Secure \& Trustworthy
Cyberspace (929), Software \& Hardware Foundations (751), Information
Integration \& Informatics (676), and Robust Intelligence (588); these are the
program names the funding autocomplete offers as its third suggestion group
(\S\ref{sec:fundingui}).

Finally, collaborative awards are the reason the interface reports project
value separately at all. The largest matched collaboration in the window is a
\$12.0M National AI Research Resource award split across 7 institutional
portions, followed by four \$10.0M Expeditions-in-Computing collaborations
spanning 4 to 10 portions each. Attributed per investigator, these become
unremarkable figures; read as projects, they are among the largest things in
the data.

\section{Limitations}
\label{sec:limitations}

\paragraph{The proxy is upstream of everything.} \sys{} inherits \csr{}'s
premise that fractional publication credit at selective conferences is a useful
proxy for research activity. Work published in journals, at workshops, or at
venues outside the curated lists is invisible, as is teaching, mentorship,
software, and service. Our stability analysis quantifies sensitivity
\emph{within} this framework; it says nothing about the framework's validity.

\paragraph{Aggregated inputs.} Because the publication table carries no author
order or paper identity, alternative credit models --- first-author,
last-author-weighted, harmonic --- cannot be evaluated at corpus scale without
per-author DBLP retrieval. We regard this as the most interesting blocked
direction.

\paragraph{Historical attribution quality.} OpenAlex affiliation histories
contain errors, which is why a community-corrections file is merged over them.
Results computed under historical attribution should be treated as
provisional.

\paragraph{Fragility is a counterfactual.} The analysis assumes a department's
remaining output is unchanged when someone leaves, which no real department
experiences: positions are refilled, collaborations are redistributed, and
students continue. It measures how concentrated current output is, expressed in
units of departures, not a forecast.

\paragraph{Name identity.} All joins are by name. \csr{} disambiguates with
numeric suffixes, which we preserve, but coauthor matching against DBLP remains
exact-match by construction: we prefer to miss a match than to attribute one
researcher's collaborators to another.

\paragraph{Funding attribution is a lower bound, and an uneven one.} Three
distinct effects bound the funding view. The equal split among listed
investigators is an assumption NSF's data cannot confirm or refute. Awards are
retrieved per current roster member and matched against their \emph{current}
affiliation, so a grant held at a previous institution, or by a departed or
retired investigator, is dropped rather than reassigned --- which makes every
departmental total a lower bound whose tightness varies with faculty turnover.
And a PI whose name resolution fails is silently absent. The Texas A\&M case in
Table~\ref{tab:nsfgap} --- 34th on publications, 167th on funding --- is large
enough that we would treat it as a matching artifact to be investigated rather
than a substantive finding, and we report it as such rather than omitting it.

\paragraph{Subfield boundaries are inherited, not derived.} The 27 top-level
areas and the venue-to-area assignment come from \csr{}. Results in
\S\ref{sec:fields} about subfield size and growth are therefore statements
about that taxonomy: a venue reassigned upstream moves output between
subfields, and work at a venue assigned to no area is counted nowhere. The
striking machine-learning growth figure in particular reflects both genuine
growth and the expansion of what the taxonomy counts as machine learning.

\paragraph{Point-in-time data.} All results are computed against upstream data
as fetched, and \csr{} updates continuously. Every number here is regenerated
by a single script, so figures can be refreshed, but any specific value in this
paper should be read as of its generation date rather than as a stable
property.

\section{Ethical considerations}
\label{sec:ethics}

Rankings applied to individuals invite misuse in ways that rankings applied to
institutions do not. We adopt several constraints. \sys{} presents no
individual-level metric as an evaluation, and reports faculty output only as the
publication counts its sources already make public. Every derived quantity is
labeled with what it measures and what it cannot, inline rather than in
documentation. The system infers no protected characteristics. Because it is
entirely client-side, it collects no record of queries.

\paragraph{Why some individuals are named and others are not.} The researcher
tables of \S\ref{sec:researchers} and the funding table of \S\ref{sec:nsf} name
people; the fragility analysis of \S\ref{sec:fragility} does not, and the line
between them is deliberate. Tables~\ref{tab:topresearchers}--\ref{tab:arealeaders}
report publication counts that \csr{} and DBLP already publish per person, in
the same units, and the tool displays them because an applicant looking for an
advisor needs exactly this. The fragility procedure produces something
different in kind: a ranked list of the individuals whose departure would most
damage their employer, which exists nowhere in the public record and which
invites personnel conclusions this work should not support. Its analytic content
--- how concentrated a department's output is --- survives fully in the
aggregate, so the names are withheld with nothing lost. The concentration
results of \S\ref{sec:concentration} are reported distributionally for the same
reason.

\paragraph{Ordering people is more hazardous than ordering departments.}
Table~\ref{tab:arealeaders} is included partly as its own counterargument: the
median subfield leader ranks 196th on the aggregate ordering and the leader in
measurement ranks 1376th, which demonstrates that a global ranking of
researchers is mostly a ranking of subfield publication volume. We report the aggregate
ordering only alongside the per-subfield one, and the interface surfaces
researchers through search and subfield context rather than as a leaderboard.

\section{Conclusion}

\sys{} began as a companion to a guide for prospective doctoral
students~\cite{phddemystify} and became an argument about presentation. The
computation underlying \csr{} is transparent and reproducible; what a single
ordered list cannot convey is how contingent its output is on choices that were
made silently. Once the ranking pipeline is fast enough to run many times ---
roughly 100\,ms in a browser, in our implementation --- that contingency
becomes directly measurable and directly displayable.

The measurement yields a finding we did not anticipate: rank stability is
sharply stratified. The top ten departments are nearly invariant while the
median department moves 17 positions across defensible settings, which means
the ranking is most reliable in the region where readers need it least.

Reporting at three levels sharpens rather than softens that conclusion. The
per-faculty ordering disagrees with the \csr{} ordering by a median of 13
positions and an NSF-funding ordering by 12, so the departmental number is one
of several defensible aggregates rather than the aggregate. Half of all
measured output now sits in four AI-adjacent subfields, so a departmental total
is dominated by a composition question the ordinal never states. And the top
decile of individuals produces 40\% of all output, so a department is a small
sample from a heavy-tailed distribution --- which is why a single departure
resolves a quarter of the fragility cases we examined.

We believe the appropriate response is not to propose a better ranking, but to
report the envelope alongside the number, and to supply the more local
information --- which subfield, which faculty, which funding --- that the
decision actually turns on.

\bibliographystyle{plain}
\bibliography{refs}

\end{document}
